% Visual style of the two-page report: Times-like serif (TeX Gyre Termes), near-black ink, no colour outside screenshots and icons.
% Compiled with XeTeX (tectonic).
\usepackage{fontspec}
\setmainfont{texgyretermes}[Extension=.otf, UprightFont=*-regular, BoldFont=*-bold, ItalicFont=*-italic, BoldItalicFont=*-bolditalic]% metric clone of Times New Roman
\setmonofont{texgyrecursor}[Extension=.otf, UprightFont=*-regular, BoldFont=*-bold, Scale=0.95]
\newcommand{\lightfont}{\mdseries}
\newcommand{\mediumfont}{\bfseries}

\usepackage[a4paper,margin=14mm]{geometry}
\usepackage{xcolor}
\usepackage{tikz}
\usetikzlibrary{arrows.meta,positioning,calc,fit,backgrounds}
\usepackage{graphicx}
\usepackage{enumitem}
\usepackage{array}
\usepackage{colortbl}
\usepackage{microtype}
\usepackage{etoolbox}
\usepackage{xstring}

\definecolor{ink}{HTML}{0A0B0D}
\definecolor{secondary}{RGB}{255,117,92}
\definecolor{hairline}{HTML}{0A0B0D}% rules and frames are black
\definecolor{grey}{HTML}{6B6E73}
\definecolor{wash}{HTML}{F3F4F5}
\color{ink}

\setlength{\parindent}{0pt}
\setlength{\parskip}{0pt}
\setlength{\columnsep}{6mm}
\pagestyle{empty}
\setlist{nosep, leftmargin=*, itemsep=0.8pt, topsep=0pt, label=\textcolor{ink}{\textbullet}}
\linespread{1.02}
\emergencystretch=2em

% Section title in sentence case; the optional argument adds space above it.
\newcommand{\rsection}[2][0mm]{%
  \par\vspace{\dimexpr 4.5mm+#1\relax}%
  {\fontsize{12}{14}\selectfont\mediumfont #2\par}%
  \vspace{1.5mm}}

% Screenshot with rounded corners; a named grey placeholder if the file is missing.
\newsavebox{\shotbox}
\newlength{\shotw}
\newlength{\shoth}
\newcommand{\shotratio}{0.52}
\newcommand{\shotzoom}{1}
\newcommand{\shotfocus}{0}
\newcommand{\shotfocusx}{0.5}% 0 shows the left edge, 0.5 the centre, 1 the right edge
\newcommand{\shot}[3]{%
  \setlength{\shotw}{#2}%
  \setlength{\shoth}{\shotratio\shotw}%
  \begin{minipage}[t]{\shotw}\centering
    \IfFileExists{images/#1}{%
      \sbox{\shotbox}{\includegraphics[width=\shotzoom\shotw]{images/#1}}%
      \ifdim\ht\shotbox<\shoth \sbox{\shotbox}{\includegraphics[height=\shoth]{images/#1}}\fi
      \begin{tikzpicture}[baseline=(current bounding box.north)]
        \useasboundingbox (0,0) rectangle (\shotw,\shoth);
        \clip[rounded corners=2mm] (0,0) rectangle (\shotw,\shoth);
        \node[anchor=north west, inner sep=0] at ({-\shotfocusx*(\wd\shotbox-\shotw)},{\shoth+\shotfocus*(\ht\shotbox-\shoth)}) {\usebox{\shotbox}};
      \end{tikzpicture}%
    }{%
      \begin{tikzpicture}[baseline=(current bounding box.north)]
        \fill[wash, rounded corners=2mm] (0,0) rectangle (\shotw,\shoth);
        \node[font=\fontsize{7}{8}\selectfont\ttfamily, text=secondary] at (0.5\shotw,0.5\shoth) {\detokenize{#1}};
      \end{tikzpicture}%
    }\\[0.8mm]
    {\fontsize{7}{8}\selectfont\color{secondary}#3}%
  \end{minipage}}

% Technical-choice table: domain and choice, reason, discarded alternative with the reason it was discarded.
% The first line of the table sits on the first baseline of the bullets next to it (both blocks are top-aligned).
\newenvironment{choixtable}{%
  \fontsize{8}{9.4}\selectfont
  \setlength{\tabcolsep}{2pt}%
  \begin{tabular*}{\linewidth}[t]{@{}>{\raggedright\arraybackslash}p{0.31\linewidth}@{\extracolsep{\fill}}>{\raggedright\arraybackslash}p{0.30\linewidth}>{\raggedright\arraybackslash}p{0.34\linewidth}@{}}%
  {\mediumfont Domaine et choix} & {\mediumfont Pourquoi} & {\mediumfont Alternative écartée}\\[0.2mm]
  \arrayrulecolor{ink}\hline\noalign{\vspace{0.4mm}}}{%
  \end{tabular*}}
\newcommand{\choix}[4]{\textbf{#1}\newline #2 & #3 & \textcolor{secondary}{#4}\\[0.2mm]\arrayrulecolor{hairline}\hline\noalign{\vspace{0.4mm}}}

% Difficulty block: a short labelled paragraph behind a thin left rule.
\newcommand{\difficulte}[3]{%
  \par\vspace{2.5mm}\noindent%
  \begin{tikzpicture}[baseline=(box.north)]
    \node[inner sep=0, text width=\dimexpr\linewidth-3mm\relax, align=justify, anchor=north west] (box) at (3mm,0) {%
      \fontsize{8.5}{10.2}\selectfont
      {\fontsize{8.5}{10.2}\selectfont\mediumfont Difficulté :}\hspace{1.5mm}#2\par\vspace{0.5mm}#3};
    \draw[hairline, line width=0.8pt] ([xshift=-1.5mm]box.north west) -- ([xshift=-1.5mm]box.south west);
  \end{tikzpicture}\par}


% Shared TikZ styles for the report figures: line styles stand for flow types, no colour.
\definecolor{lightgrey}{HTML}{B9BBBF}
\tikzset{
  lab/.style={font=\fontsize{7}{8}\selectfont, text=ink, align=center, inner sep=1pt},
  sub/.style={font=\fontsize{6}{7}\selectfont, text=secondary, align=center, inner sep=1pt},
  flowT/.style={-{Stealth[length=1.7mm,width=1.3mm]}, line width=0.6pt, draw=ink},
  flowC/.style={-{Stealth[length=1.7mm,width=1.3mm]}, line width=0.6pt, draw=ink, dash pattern=on 3pt off 2pt},
  flowV/.style={line width=0.8pt, draw=ink, dash pattern=on 0.2pt off 2.2pt, line cap=round},
  flowD/.style={-{Stealth[length=1.7mm,width=1.3mm]}, line width=0.6pt, draw=lightgrey},
  flowJ/.style={-{Stealth[length=1.5mm,width=1.1mm]}, line width=0.4pt, draw=ink},
  lane/.style={draw=hairline, line width=0.4pt, rounded corners=2mm},
}
\providecommand{\mk}[1]{node[circle, fill=ink, text=white, inner sep=0.7pt, minimum size=3.1mm, font=\fontsize{6}{6}\selectfont\bfseries, midway]{#1}}
\providecommand{\ico}[4]{%
  \node[inner sep=0] (#1) at (#2) {\includegraphics[height=7mm]{icons/#3.pdf}};
  \node[lab, below=0pt of #1] (#1l) {#4};}
\providecommand{\icoa}[4]{%
  \node[inner sep=0] (#1) at (#2) {\includegraphics[height=7mm]{icons/#3.pdf}};
  \node[lab, above=0pt of #1] (#1l) {#4};}
\providecommand{\icow}[4]{%
  \node[inner sep=0] (#1) at (#2) {\includegraphics[height=7mm]{icons/#3.pdf}};
  \node[lab, left=1pt of #1] (#1l) {#4};}

\usepackage[hidelinks]{hyperref}
